\documentclass[11pt,a4paper]{article}
\usepackage[spanish,es-noquoting]{babel}
\usepackage[top=2.2cm, bottom=2.2cm, left=2.5cm, right=2.5cm]{geometry}
\usepackage{fontspec}
\usepackage{amsmath}
\usepackage[version=4]{mhchem}
\usepackage{enumitem}
\usepackage{microtype}

\setmainfont{Calibri}
\linespread{1.15}

\begin{document}
\pagestyle{empty}

\noindent\textbf{1.} En un reactor industrial se introducen $140\text{ g}$ de gas nitrógeno (\ce{N2}) y $36\text{ g}$ de gas hidrógeno (\ce{H2}) para sintetizar amoníaco gaseoso (\ce{NH3}) según la siguiente reacción química (sin ajustar):
\[
\ce{N2 (g) + H2 (g) -> NH3 (g)}
\]

\begin{enumerate}[label=\textbf{\alph*)}, leftmargin=2em, itemsep=4pt]
    \item Escribe y ajusta la ecuación química correspondiente a la síntesis del amoníaco.
    \item Determina de forma justificada cuál de los dos reactivos actúa como reactivo limitante y cuál en exceso.
    \item Calcula la masa teórica máxima de amoníaco (\ce{NH3}) que se obtendría si la reacción fuese completa ($100\%$ de rendimiento).
    \item Sabiendo que el rendimiento real del proceso es del $80\%$, calcula la masa de amoníaco que se obtiene en la práctica.
    \item Determina los moles y los gramos de reactivo en exceso que quedan sin reaccionar.
\end{enumerate}

\vspace{0.25cm}
\noindent\textbf{Datos:} Masas atómicas ($\text{u}$): $\mathrm{H} = 1$; $\mathrm{N} = 14$.

\vspace{0.6cm}

\noindent\textbf{2.} En un horno de laboratorio se calientan $250\text{ g}$ de carbonato de calcio (\ce{CaCO3}) puro, descomponiéndose térmicamente en óxido de calcio (\ce{CaO}) y dióxido de carbono gaseoso (\ce{CO2}) según la ecuación (sin ajustar):
\[
\ce{CaCO3 (s) -> CaO (s) + CO2 (g)}
\]

\begin{enumerate}[label=\textbf{\alph*)}, leftmargin=2em, itemsep=4pt]
    \item Escribe y ajusta la ecuación química de la descomposición térmica.
    \item Calcula la masa teórica de óxido de calcio (\ce{CaO}) que se obtendría con un rendimiento del $100\%$.
    \item Si al finalizar la reacción se han obtenido experimentalmente $112\text{ g}$ de óxido de calcio, calcula el porcentaje de rendimiento ($\%$) del proceso.
    \item Calcula la masa real de dióxido de carbono (\ce{CO2}) que se ha desprendido durante el experimento.
\end{enumerate}

\vspace{0.25cm}
\noindent\textbf{Datos:} Masas atómicas ($\text{u}$): $\mathrm{C} = 12$; $\mathrm{O} = 16$; $\mathrm{Ca} = 40$.

\end{document}
